    \caption{\ai{\textbf{Pre-stimulus state and single-trial variability in the impulse experiments.}}
    \ai{State measures as in Fig.~\ref{fig:figure2}G, each over the 1\,s before onset; 1516 trials in 4 sessions from 3 mice (Mouse 5a and 5b are two sessions of one mouse).}
    \ai{\textbf{A--C:} What relative and absolute 2--4\,Hz power measure.}
    \ai{\textbf{A:} Relative versus absolute 2--4\,Hz power for every trial, each $z$-scored within session (absolute power on a log scale); Spearman $\rho = 0.47$. Circles mark the trials in \textbf{B}.}
    \ai{\textbf{B:} Three trials from Mouse 1b chosen by rule, on one scale: typical (both measures near the session mean), high absolute and low relative power ($z_\mathrm{abs} = +0.9$, $z_\mathrm{rel} = -1.5$), and the highest relative power among trials with at least average absolute power ($z_\mathrm{abs} = +1.1$, $z_\mathrm{rel} = +2.2$). The last two have similar absolute 2--4\,Hz power. Gray shading, the window over which power was measured; red, laser command.}
    \ai{\textbf{C:} Mean spectra over the same window for Mouse 1b trials that were typical ($|z| < 0.5$ on both; $n = 72$), high in absolute and below average in relative power ($z_\mathrm{abs} > 1$, $z_\mathrm{rel} < 0$; $n = 20$), or high in relative power with at least average absolute power ($z_\mathrm{rel} > 1$, $z_\mathrm{abs} \ge 0$; $n = 104$). Left, power spectral density; right, each frequency bin as a fraction of 0.4--10\,Hz power. Gray band, 2--4\,Hz. Over all trials, absolute 2--4\,Hz power followed total 0.4--10\,Hz power (Spearman $\rho = 0.97$), whereas relative power did not ($\rho = 0.24$).}
    \ai{\textbf{D--I:} Variability by state quartile, session by session. Rows, motion energy (\textbf{D}, \textbf{G}), relative 2--4\,Hz power (\textbf{E}, \textbf{H}) and absolute 2--4\,Hz power (\textbf{F}, \textbf{I}). Colors, sessions.}
    \ai{\textbf{D--F:} From left, root-mean-square trial-to-trial variability in each quartile, normalized to the session's typical trial (dotted line, 1; red, mean of sessions; gray dashed, the stimulus-free control); response variability divided by stimulus-free variability, per quartile, which stays flat for a state that acts only on ongoing activity (black, mean of sessions); and single-trial errors per quartile and session (box, 25th--75th percentile; whiskers, 10th--90th; line, median). Annotations: top versus bottom quartile, geometric mean of the four per-session ratios with their range.}
    \ai{\textbf{G--I:} Single trials at one stimulation amplitude per session, chosen by a fixed rule (each session's second-highest amplitude: Mouse 5a, 0.80\,mW, 59 trials; Mouse 5b, 0.87\,mW, 50; Mouse 1b, 1.43\,mW, 82; Mouse 3c, 0.33\,mW, 100), one trial per row, sorted by the row's state with the lowest at the bottom; horizontal lines, quartile boundaries. Activity, $\Delta F/F$ from 0.2\,s before to 0.49\,s after onset (dashed line), blue negative to red positive over $\pm 5\%$. $|$Dev.$|$, the absolute difference between each trial and the mean of all trials at that amplitude, 0--0.49\,s after onset, white (0) to black (5\% $\Delta F/F$). Both in raw $\Delta F/F$, not normalized. Descriptive figure; supporting tests are in Table~\ref{tab:stat_fig2g}.}}
